\section{Teaching, Asceticism, and the Making of a School}
\label{sec:origen-school}

\subsection{An eighteen-year-old catechist}

During continuing danger at Alexandria, Origen began teaching people who came
for Christian instruction. Eusebius places him in charge around age eighteen
and eventually says Bishop Demetrius entrusted the work to him. Plutarch,
Heraclas, two men named Serenus, Heraclides, Hero, Herais, and others are
remembered as pupils who suffered in persecution (\emph{Ecclesiastical
History} 6.3--5). The list makes the classroom part of martyr history: learning
was preparation for baptism and possible public confession, not an insulated
academic career.

The institutional language must remain proportionate. Origen probably taught
within an episcopally supervised Alexandrian catechetical structure, but the
form, curriculum, and boundary between private school, church instruction, and
advanced seminar changed over time. Eusebius says Origen eventually divided
the work: Heraclas taught beginners while Origen concentrated on advanced
students and Scripture (6.15). Emanuela Prinzivalli reads this division within
Demetrius's effort to bring Alexandrian Christian teaching under more regular
episcopal direction. That is a named modern reconstruction, not a surviving
constitution.

Pupils did not all begin inside the Church. Philosophically educated pagans and
members of rival Christian schools heard him. Origen's own defense, quoted by
Eusebius at 6.19.12--14, says he studied philosophical systems because learned
and heterodox interlocutors came to him. He names Pantaenus and Heraclas as
precedents. Philosophy in this account is both intellectual equipment and
missionary risk. It could clarify an opponent's premises; it could also furnish
categories later critics thought had governed Scripture.

\subsection{A severe experiment in Christian philosophy}

Origen sold literary books for a small daily income, reduced food and sleep,
went barefoot for years, and rejected offers of support, according to Eusebius
(6.3). Such details belong to ancient philosophical biography as well as
Christian ascetic memory. Their literary fit does not make them false, but it
warns against treating every austerity as a measured diary entry.

The most consequential report says the young teacher took Matthew 19:12
physically and castrated himself. Eusebius presents the act as immature zeal,
partly intended to prevent scandal because Origen instructed women, and says
Demetrius later used it against his ordination (6.8). Origen's later
\emph{Commentary on Matthew} 15.1--5, extant in Greek, rejects genital
mutilation as the meaning of the verse and redirects the command toward moral
continence. Modern scholars remain divided. Henri Crouzel and Joseph Trigg
accept Eusebius's report while treating the commentary as later correction;
R. P. C. Hanson likewise defended its historicity. Henry Chadwick, Patricia
Cox, and Mariusz Szram give skeptical or deliberately unresolved readings,
while Daniel Caner establishes that the bodily practice was historically
intelligible without thereby supplying corroboration. The later rejection can
represent maturation, but the story also served ecclesial controversy and fits
hostile caricature.
No subject-authored admission, medical record, or contemporary independent
witness survives. This biography therefore records the act as an early,
specific, disputed report---neither a secure bodily fact nor a legend safe to
discard.

Whatever happened physically, Origen's ascetic program was real. Celibacy,
fasting, poverty, wakefulness, study, and willingness to suffer made the
teacher's body part of his argument. Gregory's later farewell address praises
the consistency of life and teaching. The cost must remain visible: Eusebius
himself says the regime threatened Origen's health. Zeal is historically
explanatory without making bodily damage exemplary.

\subsection{Learning Greek philosophy without becoming a Greek school}

Porphyry, the anti-Christian philosopher, supplied an invaluable hostile
witness later quoted by Eusebius. He names Origen's study of Plato, Numenius,
Cronius, Pythagorean writers, Chaeremon, and Cornutus, and mocks his use of
Greek learning to interpret Jewish Scriptures (6.19). Eusebius disputes
Porphyry's account of Origen's Christian origins but accepts his erudition.
Their convergence establishes serious philosophical study; it does not settle
which Ammonius taught Origen or permit reconstruction of a lost syllabus.

Identifying Origen's teacher with the Ammonius Saccas later associated with
Plotinus is possible but not proved. Nor is ``Platonist'' a sufficient summary
of Origen. He borrowed philosophical vocabulary and argumentative forms while
opposing an uncreated or coequal cosmos or matter, fixed classes of souls,
polytheistic cult, and other positions he judged incompatible with Scripture.
Whether eternal divine creativity implied successive ages was a different,
disputed question in his works and their transmission. In the
\emph{Letter to Gregory}, philosophy and the mathematical disciplines become
preparatory instruments; Scripture remains the end, and prayer remains
necessary to open it.

The result was not Christianity translated without remainder into Platonism.
It was a school in which the Bible, the Church's proclamation, ancient
philology, philosophical dispute, moral formation, and prayer interrogated one
another. Later reception repeatedly broke that unstable unity apart, keeping
the exegete while prosecuting the speculative philosopher.

\begin{studyframe}
\textbf{Ascetic evidence.} Eusebius's Origen is a Christian philosopher whose
life authenticates his teaching. That is evidence for an ancient moral ideal
and for remembered practices. It is not permission to diagnose Origen, certify
the disputed operation, or recommend his austerities without pastoral and
medical judgment.
\end{studyframe}
